\chapter{Посадка в клетку}\label{ch:posadka}

\textbf{Что вы узнете:} почему дрон садится не точно в центр клетки и как это исправить командой \kod{land_on}.

\section{land и land\_on}

Команда \kod{d.land()} сажает дрона там, где он оказался. Но позиция дрона точна лишь до 10--15\cm, и после посадки он оказывается рядом с клеткой, а не в её центре. Если вы хотите, чтобы дрон вернулся ровно на свою клетку (чтобы потом снова стартовать с неё), используйте \kod{d.land_on(cell)}.

\solfile{l08_landon.py}

Дрон облетает квадрат и садится в центр клетки 13.

\section{Что делает land\_on}

Команда \kod{land_on(cell, seconds=2, low=35)} — это четыре шага подряд:
\begin{enumerate}
  \item \kod{code(cell, seconds)}: подлететь к клетке (по умолчанию 2\,с);
  \item \kod{set_z(z, 1.5)}: повисеть 1,5\,с, чтобы успокоиться;
  \item \kod{set_z(low, 2.5)}: плавно снизиться до \kod{low} (по умолчанию 35\cm) за 2,5\,с;
  \item \kod{land()}: сесть (3\,с).
\end{enumerate}
Здесь \kod{z} — высота, на которой дрон был до этого. Если взлёта в программе ещё не было, будет ошибка «land\_on: сначала нужен взлёт».

\begin{figure}[H]
  \centering
  % land_on: высота по времени
\begin{tikzpicture}[x=0.75cm, y=0.032cm, font=\sffamily\small]
  \draw[-{Stealth}, ink] (-0.2,0) -- (11.6,0) node[right] {$t$, с};
  \draw[-{Stealth}, ink] (0,-4) -- (0,125) node[above] {$z$, см};
  \draw[accent, very thick] (0,100) -- (3.5,100) -- (6,35) -- (9,0);
  \draw[dashed, rule] (2,0) -- (2,100); \draw[dashed, rule] (3.5,0) -- (3.5,100); \draw[dashed, rule] (6,0) -- (6,35);
  \node[font=\scriptsize, text=inkgray, anchor=south] at (1,103) {code(cell, 2)};
  \node[font=\scriptsize, text=inkgray, anchor=south] at (2.75,103) {set\_z};
  \node[font=\scriptsize, text=inkgray, anchor=south west] at (5.3,82) {set\_z(35, 2.5)};
  \node[font=\scriptsize, text=inkgray, anchor=south west] at (7.3,22) {land()};
  \foreach \t in {0,2,3.5,6,9}{\draw[ink] (\t,-2) -- (\t,2); \node[below, font=\scriptsize] at (\t,-3) {\t};}
  \foreach \z in {35,100}{\draw[ink] (-0.07,\z) -- (0.07,\z); \node[left, font=\scriptsize] at (-0.1,\z) {\z};}
\end{tikzpicture}

  \caption{Высота дрона при \texttt{land\_on(cell)} с 100\cm: подлёт 2\,с, зависание 1,5\,с, снижение до 35\cm за 2,5\,с, посадка 3\,с. Всего 9 секунд.}\label{fig:landon}
\end{figure}

\begin{zametka}
\kod{land_on} стоит использовать, когда дрон должен вернуться на свою клетку или сесть в заданную. Он делает шоу на 6 секунд длиннее, чем простое \kod{land()} (9 вместо 3).
\end{zametka}

\begin{probuy}
  \item Из клетки 13 в клетку 16 и посадка в центр клетки 16.
  \item То же, но \kod{land_on(16, 3, 60)}: подлёт 3 секунды, остановка на 60\cm. Как изменится длительность шоу?
  \item Замените \kod{land_on(16)} четырьмя обычными командами с теми же временами (из списка выше).
  \item Из клетки 13 в клетку 52 по диагонали и посадка в центр клетки 52.
\end{probuy}
\otvety{8}
